\documentclass[11pt]{article}
\usepackage[a4paper,margin=27mm]{geometry}
\usepackage[T1]{fontenc}
\usepackage[utf8]{inputenc}
\usepackage{lmodern}
\usepackage{amsmath,amssymb,amsthm,mathtools}
\usepackage{booktabs}
\usepackage{xcolor}
\usepackage{microtype}
\usepackage{hyperref}
\hypersetup{colorlinks=true,linkcolor=blue!45!black,citecolor=blue!45!black,urlcolor=blue!55!black}

\newtheorem{theorem}{Theorem}[section]
\newtheorem{corollary}[theorem]{Corollary}
\newtheorem{lemma}[theorem]{Lemma}
\newtheorem{remark}[theorem]{Remark}
\newcommand{\Res}{\operatorname{Res}}
\newcommand{\Disc}{\operatorname{Disc}}
\newcommand{\ord}{\operatorname{ord}}

\title{\textbf{Cyclotomic Product-Row Support:}\\
Principal Subresultants, Fitting Loci, and Resultant Transversality}
\author{Maksim Churilov}
\date{Research companion --- 5 September 2026}

\begin{document}
\maketitle

\begin{abstract}
We sharpen an exact arithmetic support problem arising from fourth-moment
reductions for arithmetic quantum cat maps. For odd prime order $K=2m+1$, a
fixed-product decomposition reduces the non-forced Cayley-energy excess to
intersections of special affine copies of the prime cyclotomic trace set. A
specialization-safe maximal-real formulation associates to every normalized
row ratio $r\ne1$ two degree-$m$ polynomials $P,Q$ over
$\mathbf Z[\tau]$, $\tau=\zeta_K+\zeta_K^{-1}$, such that over every split
residue embedding the row-intersection number equals $\deg\gcd(P,Q)$.

We prove that the order-three event is detected by only the first three
principal subresultant coefficients
\[
\operatorname{psc}_0(P,Q),\quad
\operatorname{psc}_1(P,Q),\quad
\operatorname{psc}_2(P,Q),
\]
and identify the same event intrinsically as the third-corank Fitting locus of
the $2m\times2m$ Sylvester presentation. Thus the previously used full
subresultant coefficient ideal and scalar shift-resultant coordinates are
different coordinate systems on one determinantal stratum.

For the natural row translation
\[
\mathcal R_r(Z)=\Res(P,Q_Z)=c_0+c_1Z+c_2Z^2+\cdots,
\]
we prove a transversality formula at every stratum. If
$P=GA$, $Q=GB$, $d=\deg G$, then the first nonzero shift coefficient satisfies,
up to the fixed resultant/subresultant sign convention,
\[
c_d=\alpha_r^{-d}\operatorname{psc}_d(P,Q)
\Disc(G)\Res(A,G)\Res(G,B).
\]
Every extra factor is nonzero on the maintained squarefree split locus.
Hence the shift and principal-subresultant coordinates are transverse local
coordinates on the same gcd stratification.

For the inverse ratio $r=-1$, reflection halves the degree. Writing
$F(X)=E(X^2)+XO(X^2)$, the inverse shift resultant factors
\[
\Res_X(F(X),F(Z-X))=2^mF(Z/2)S(Z)^2,
\quad
S(Z)=\prod_{i<j}(Z-x_i-x_j),
\]
and furthermore
\[
S(Z)=\pm\Res_Y(E_Z,O_Z),
\quad F(X+Z/2)=E_Z(X^2)+XO_Z(X^2).
\]
At a unique common root $y$ of $E,O$ the translation derivative is
\[
\rho'(0)=\pm\Res\!\left(\frac E{Y-y},\frac O{Y-y}\right)
\bigl(E'(y)^2-yO'(y)^2\bigr),
\]
where the final factor equals
$-F'(z)F'(-z)/(4y)$ for $z^2=y$. Squarefreeness therefore forces
transversality at gcd degree one; vanishing of both $\rho(0)$ and $\rho'(0)$
is exactly the degree-at-least-two event.

We further localize the inverse event to four roots: $I_{-1}\ge3$ if and only
if four distinct centred roots have $e_1=e_3=0$. The pair-sum polynomial obeys
the universal factorization
\[
\Disc(S_F)=\Disc(F)^{m-2}\prod_{|A|=4}
\bigl(e_1(A)^3-4e_1(A)e_2(A)+8e_3(A)\bigr)^2.
\]
Writing $v_n=(t_n-2)/(\tau-2)$ and
$w_{a,b}=2-v_a-v_b$ removes the universal ramified conductor factor:
$2\tau-t_a-t_b=(\tau-2)w_{a,b}$ and $|N(\tau-2)|=K$.
Thus inverse order-three support is a union of two-generator local ideal
conditions $(w_e,w_f)$ for disjoint edges. An exact finite certificate at
$(K,p)=(89,143207231)$ has $I_{-1}=4$ and
$p=2.282469\ldots K^4$, disproving a strict finite $p\le K^4$ ceiling while
remaining compatible with the open asymptotic target $K^{4+o(1)}$.

These results reduce the remaining pointwise arithmetic problem to common
split prime divisors of small canonical cyclotomic systems. They do not prove a
polynomial largest-prime bound and do not prove microscopic BGS universality.
\end{abstract}

\section{Fixed-degree row formulation}

Let $K=2m+1\ge7$ be an odd prime, let
$\zeta=\exp(2\pi i/K)$, and put $\tau=\zeta+\zeta^{-1}$. The full subgroup
trace polynomial is
\[
H_K(T)=(T-2)C_K(T),
\]
where $C_K$ is the minimal polynomial of $\tau$. For a normalized row ratio
$r\ne1$ the affine trace map has the Chebyshev form
\[
Y=\alpha_rX+\beta_r,
\qquad
\alpha_r=U_{r-1}(\tau/2),
\qquad
\beta_r=-2U_{r-2}(\tau/2).
\]
The two full row trace polynomials share one characteristic-zero trace
corresponding to the invalid row class. Removing this trace before reduction
gives fixed-degree cofactors
\[
P_{K,r}(T),Q_{K,r}(T)\in\mathbf Z[\tau][T],
\qquad
\deg P_{K,r}=\deg Q_{K,r}=m.
\]
At every rational prime $p\equiv1\pmod K$ and every embedding prime
$\mathfrak q\mid p$, the leading coefficients remain nonzero and
\begin{equation}\label{eq:gcd-count}
I_{1,r}(\mathfrak q)
 =\deg\gcd(P_{K,r},Q_{K,r}\bmod\mathfrak q).
\end{equation}
This fixed-degree statement is the specialization safeguard. Earlier
projective dehomogenizations can lose degree at the point at infinity and are
therefore retained only as one-way obstructions.

\section{Principal subresultants are sufficient}

For two fixed-degree polynomials $P,Q$ over a field, let $S_j(P,Q;T)$ be the
standard $j$th subresultant polynomial and let
\[
\operatorname{psc}_j(P,Q)=[T^j]S_j(P,Q;T).
\]
We use the standard subresultant theorem \cite{BrownTraub}: if
$g=\deg\gcd(P,Q)$, then
\[
S_0=\cdots=S_{g-1}=0,
\qquad
S_g=c\gcd(P,Q),\qquad c\ne0.
\]

\begin{lemma}[Principal coefficients detect gcd degree]\label{lem:psc}
For every $d\ge1$,
\[
\deg\gcd(P,Q)\ge d
\iff
\operatorname{psc}_0(P,Q)=\cdots=
\operatorname{psc}_{d-1}(P,Q)=0.
\]
\end{lemma}

\begin{proof}
If the gcd degree is at least $d$, the first $d$ subresultants vanish. If the
gcd degree is $g<d$, the first nonzero subresultant is
$c\gcd(P,Q)$ of degree exactly $g$, so its $T^g$ coefficient is nonzero.
\end{proof}

Apply this with $d=3$ to the fixed-degree row cofactors. Define
\[
\pi_j=\operatorname{psc}_j(P_{K,r},Q_{K,r}),
\qquad j=0,1,2,
\]
and
\[
J^{\mathrm{psc}}_{K,r}=(\pi_0,\pi_1,\pi_2)
 \subset\mathbf Z[\tau].
\]
Let
\[
J^{\mathrm{psc}}_{K,r}\cap\mathbf Z
 =\Pi_{K,r}^{(3)}\mathbf Z.
\]

\begin{theorem}[Three-scalar exact split support]\label{thm:psc-support}
For every split rational prime $p\equiv1\pmod K$,
\[
p\mid\Pi_{K,r}^{(3)}
\iff
\text{some }\mathfrak q\mid p\text{ has }I_{1,r}(\mathfrak q)\ge3.
\]
\end{theorem}

\begin{proof}
Combine \eqref{eq:gcd-count} with Lemma~\ref{lem:psc}. A rational prime divides
the contraction of an ideal precisely when the ideal is contained in some
prime ideal above it.
\end{proof}

Thus the large full-coefficient subresultant support and the three principal
coefficients have the same split rational-prime support. Their valuations need
not agree. Exact Smith-normal-form controls include
\[
\begin{array}{c|c|c}
(K,r)&\Pi_{K,r}^{(3)}&\text{prime-to-}2K\text{ part}\\\hline
(7,2)&7&1\\
(11,7)&11^2&1\\
(13,3)&13^3&1\\
(13,12)&2^4\,5^3\,13^3&5^3\\
(17,10)&2^2\,17^5\,103&103.
\end{array}
\]
At $(K,r)=(13,12)$ the independently maintained full/shift support has
prime-to-$2K$ valuation $5^2$ rather than $5^3$. This is an explicit negative
control against identifying the support integers themselves.

\section{The intrinsic Sylvester/Fitting locus}

Let
\[
\mathcal S_{P,Q}:k[T]_{<m}\oplus k[T]_{<m}
 \longrightarrow k[T]_{<2m},
\qquad
(U,V)\mapsto UP+VQ.
\]
In monomial bases this is the $2m\times2m$ Sylvester matrix.

\begin{theorem}[Sylvester nullity]\label{thm:nullity}
For same-degree nonzero polynomials over a field,
\[
\dim\ker\mathcal S_{P,Q}=\deg\gcd(P,Q).
\]
\end{theorem}

\begin{proof}
Write $P=GP_1$, $Q=GQ_1$ with $(P_1,Q_1)=1$ and $g=\deg G$. A kernel vector
satisfies $UP_1=-VQ_1$, so $U=WQ_1$, $V=-WP_1$. The degree bounds are exactly
$\deg W<g$.
\end{proof}

Consequently
\[
\deg\gcd(P,Q)\ge3
\iff
\operatorname{rank}\mathcal S_{P,Q}\le2m-3.
\]
Let $I_{2m-2}(\mathcal S)$ be the ideal generated by all
$(2m-2)\times(2m-2)$ minors. Then
\[
\deg\gcd(P,Q)\ge3
\iff I_{2m-2}(\mathcal S)=0
\]
after specialization to a field. Equivalently, the
$(2m-2)$nd compound matrix vanishes. If the gcd degree is exactly two, this
compound matrix has rank one.

Thus the order-three row support is intrinsically the third-corank Fitting
locus of a cyclotomic Sylvester presentation. The PSC triple is a small
canonical coordinate system cutting out the same field-valued support, while
the all-minor Fitting ideal exposes exterior-power and determinantal identities.

\section{Shift-resultant transversality}

Let $k$ be a field, let $f,g$ be monic degree-$n$ polynomials, and suppose
\[
f=GA,\qquad g=GB,
\]
where $G$ is monic squarefree of degree $d$ and $(A,B)=1$. Consider a
degree-preserving deformation
\[
g_Z=g+Z\psi+O(Z^2)
\]
and assume $\psi(y)\ne0$ at every root $y$ of $G$. Put
\[
R(Z)=\Res(f,g_Z).
\]

\begin{theorem}[First nonzero shift coefficient]\label{thm:shift-psc}
One has $\ord_{Z=0}R=d$, and, up to the fixed resultant/subresultant sign
convention,
\[
[Z^d]R(Z)
 =\operatorname{psc}_d(f,g)\,
  \Res(A,G)\prod_{G(y)=0}\psi(y).
\]
\end{theorem}

\begin{proof}
The roots of $f$ coming from $G$ each contribute one factor
$Z\psi(y)+O(Z^2)$. The remaining roots, those of $A$, contribute
\[
\Res(A,g)=\Res(A,G)\Res(A,B).
\]
For monic $f,g,G$ the first nonzero subresultant is
$\pm\Res(A,B)G$, so its principal coefficient is
$\pm\Res(A,B)$.
\end{proof}

For the row shift, normalize the target polynomial to be monic and write
\[
q_Z(T)=\alpha_r^{-m}
 R_r(\alpha_rT+\beta_r+Z).
\]
Then
\[
\psi(T)=\alpha_r^{-m}R_r'(\alpha_rT+\beta_r),
\qquad
q_0'(y)=\alpha_r\psi(y).
\]
On the maintained split locus the target cofactor is squarefree and
$\alpha_r\ne0$, so every $\psi(y)$ is nonzero.

If $G=\gcd(P,q_0)$, the product in
Theorem~\ref{thm:shift-psc} can be evaluated further:
\[
\prod_{G(y)=0}\psi(y)
 =\pm\alpha_r^{-d}\Disc(G)\Res(G,B).
\]
Hence
\begin{equation}\label{eq:explicit-local}
[Z^d]\Res(P,q_Z)
 =\pm\alpha_r^{-d}
 \operatorname{psc}_d(P,q_0)
 \Disc(G)\Res(A,G)\Res(G,B).
\end{equation}
Every multiplier relating the first nonzero shift coordinate to the first
nonzero PSC is nonzero in the residue field. Thus the shift coefficients and
PSCs are transverse coordinate systems on the same gcd stratification.

In particular,
\[
\begin{array}{c|c|c}
\deg\gcd&\text{PSC coordinates}&\text{shift coordinates}\\\hline
0&\pi_0\ne0&c_0\ne0\\
1&\pi_0=0,\ \pi_1\ne0&c_0=0,\ c_1\ne0\\
2&\pi_0=\pi_1=0,\ \pi_2\ne0&c_0=c_1=0,\ c_2\ne0\\
\ge3&\pi_0=\pi_1=\pi_2=0&c_0=c_1=c_2=0.
\end{array}
\]

\section{Inverse pair sums and half-degree support}

For $r=-1$, reflection gives
\[
Q(T)=-P(2\tau-T).
\]
Set
\[
F(X)=P(\tau+X)=E(X^2)+XO(X^2).
\]
Then over every split embedding prime
\[
I_{-1}=2\deg\gcd(E,O).
\]
Let $x_1,\ldots,x_m$ be the roots of $F$ and put
\[
S(Z)=\prod_{i<j}(Z-x_i-x_j).
\]

\begin{theorem}[Inverse pair-sum square]\label{thm:pairsum}
\[
\Res_X(F(X),F(Z-X))=2^mF(Z/2)S(Z)^2.
\]
At split primes $2^mF(0)$ is a unit, hence
\[
I_{-1}=2\ord_{Z=0}S,
\qquad
I_{-1}\ge3\iff S(0)=S'(0)=0.
\]
\end{theorem}

The pair-sum polynomial itself is a half-degree resultant. If
\[
F(X+Z/2)=E_Z(X^2)+XO_Z(X^2),
\]
then
\begin{equation}\label{eq:translated}
S(Z)=\pm\Res_Y(E_Z,O_Z).
\end{equation}
To first order,
\[
\partial_ZE_Z|_0=\frac{O+2YO'}2,
\qquad
\partial_ZO_Z|_0=E'.
\]
Thus inverse order-three support is a two-scalar problem: a half-degree
resultant and its translation derivative.

\section{Inverse resultant transversality}

We record a local resultant derivative identity in a form independent of the
cyclotomic application.

\begin{theorem}[Simple-common-root derivative]\label{thm:derivative}
Let
\[
f=(Y-y)A,\qquad g=(Y-y)B,
\qquad (A,B)=1,
\]
and assume $y$ is simple in both $f$ and $g$. For
$f_t=f+t\phi$, $g_t=g+t\psi$,
\[
\left.\frac d{dt}\Res(f_t,g_t)\right|_{t=0}
 =(-1)^{\deg f-1}\Res(A,B)
 \bigl(f'(y)\psi(y)-g'(y)\phi(y)\bigr).
\]
\end{theorem}

\begin{proof}
Let $y(t)$ be the perturbed root of $f_t$ with $y(0)=y$. Then
$y'(0)=-\phi(y)/f'(y)$ and
\[
g_t(y(t))
 =t\left(\psi(y)-\frac{g'(y)\phi(y)}{f'(y)}\right)+O(t^2).
\]
The other roots of $f$ contribute the nonzero product
$\Res(A,g)=\Res(A,Y-y)\Res(A,B)$, and
$\Res(A,Y-y)=(-1)^{\deg f-1}f'(y)$.
\end{proof}

Apply this to
\[
\rho(Z)=\Res(E_Z,O_Z).
\]
At a unique common root $y$ of $E,O$,
\[
\rho'(0)
 =\pm\Res\!\left(\frac E{Y-y},\frac O{Y-y}\right)
 \bigl(E'(y)^2-yO'(y)^2\bigr).
\]
If $z^2=y$, then direct differentiation of
$F(X)=E(X^2)+XO(X^2)$ gives
\[
E'(y)^2-yO'(y)^2=-\frac{F'(z)F'(-z)}{4y}.
\]
The maintained split trace polynomial is squarefree and $F(0)$ is a unit, so
this factor is nonzero whenever the gcd degree is one. Therefore
\[
\deg\gcd(E,O)=1
\Longrightarrow
\rho(0)=0,\quad \rho'(0)\ne0,
\]
while
\[
\deg\gcd(E,O)\ge2
\iff
\rho(0)=\rho'(0)=0.
\]
The second scalar is therefore a genuine transversality test, not an accidental
coefficient.

More generally, if a squarefree gcd
$G=\prod_{j=1}^d(Y-y_j)$ is shared by $f=GA$, $g=GB$, then for a directional
perturbation $f_t=f+t\phi$, $g_t=g+t\psi$ the first nonzero resultant
coefficient is
\[
[t^d]\Res(f_t,g_t)
 =\Res(A,G)\Res(A,B)
 \prod_{j=1}^d
 \frac{f'(y_j)\psi(y_j)-g'(y_j)\phi(y_j)}{f'(y_j)}.
\]
When all local Jacobians are nonzero, the vanishing order is exactly $d$.

\section{Quartic localization and reduced inverse edge ideals}

The pair-sum square admits a further localization which is independent of the
half-degree determinant coordinates.

\begin{theorem}[Four-root inverse certificate]\label{thm:quartic}
On the maintained odd, split, squarefree locus,
\[
I_{-1}\ge3
\iff I_{-1}\ge4
\iff
\text{some four distinct centred roots }x_1,\ldots,x_4
\text{ satisfy } e_1=e_3=0.
\]
\end{theorem}

\begin{proof}
By Theorem~\ref{thm:pairsum}, $I_{-1}\ge3$ means that at least two distinct
pair-sum factors vanish at zero. They cannot share a vertex, since
$x_i+x_j=x_i+x_k=0$ and squarefreeness would force $j=k$. Hence there are four
distinct roots forming two opposite pairs, which gives $e_1=e_3=0$.

Conversely, the monic quartic on four distinct roots is
\[
X^4-e_1X^3+e_2X^2-e_3X+e_4.
\]
If $e_1=e_3=0$ it is even. The roots are nonzero because $F(0)$ is a unit, so
in odd characteristic they form two distinct opposite pairs. Thus two
pair-sum factors vanish and $I_{-1}=2\ord_0S\ge4$.
\end{proof}

There is also a universal discriminant identity. For a four-subset $A$ of the
roots put
\[
q_A=e_1(A)^3-4e_1(A)e_2(A)+8e_3(A).
\]
For any labeling $A=\{i,j,k,l\}$,
\[
q_A=(x_i+x_j-x_k-x_l)(x_i+x_k-x_j-x_l)
(x_i+x_l-x_j-x_k).
\]

\begin{theorem}[Pair-sum discriminant]\label{thm:pairsum-disc}
For monic $F$ of degree $m$,
\[
\Disc(S_F)=\Disc(F)^{m-2}\prod_{|A|=4}q_A^2.
\]
\end{theorem}

\begin{proof}
Index roots of $S_F$ by edges of the complete graph on the roots of $F$.
Two edges sharing a vertex differ by one root difference. Each unordered root
pair appears in this way for exactly $m-2$ choices of shared vertex, giving
$\Disc(F)^{m-2}$. Disjoint edge pairs live on one four-subset, and its three
partitions into two edges give the three factors in $q_A$. Squaring the root
differences gives the displayed identity.
\end{proof}

For the cyclotomic trace set define
\[
v_0=0,\qquad v_n=\frac{t_n-2}{\tau-2}\quad(n\ne0),
\qquad w_{a,b}=2-v_a-v_b.
\]
Since
\[
t_n-2=s_{n/2}^2,\qquad \tau-2=s_{1/2}^2,
\]
we have
\[
v_n=\left(\frac{s_{n/2}}{s_{1/2}}\right)^2,
\]
the square of a real cyclotomic unit. Moreover
\begin{equation}\label{eq:reduced-edge}
2\tau-t_a-t_b=(\tau-2)w_{a,b},
\qquad |N_{\mathbf Q(\tau)/\mathbf Q}(\tau-2)|=K.
\end{equation}
Thus for split $p\ne K$, the inverse zero-edge condition is simply
$w_{a,b}=0\pmod{\mathfrak q}$.

For two disjoint valid edges $e,f$ put
\[
\mathfrak j_{e,f}=(w_e,w_f)\subset\mathbf Z[\tau].
\]
Theorem~\ref{thm:quartic} and primality of $\mathfrak q$ imply that the inverse
order-three support is exactly the union of the prime-ideal supports of these
two-generator ideals. Equivalently one may use the pairing-free quartic ideals
$(e_1,e_3)$ and take their product over four-subsets.

A rational gcd of norms is only a candidate filter: the two norm divisibilities
may occur at different primes above the same rational $p$. The finite value
$(K,p)=(41,2908049)$ is a maintained negative control for this failure mode.

The strongest current inverse witness is
\[
K=89,\qquad p=143207231,
\qquad \frac{p}{K^4}=2.282469174\ldots.
\]
At one split embedding there are exactly two zero pair-sum edges,
\[
(16,35),\qquad(20,33),
\]
so $I_{-1}=4$. The associated four roots satisfy $e_1=e_3=0$, and exact
cyclotomic norm computation gives
\[
|N(w_{16,35})|=|N(w_{20,33})|=143207231.
\]
An independent C++20 program verifies primality, exact order $89$, distinct
trace classes, the two zero edges and the quartic certificate directly in
$\mathbf F_p$.

This disproves the stronger finite assertion that every inverse supported
split prime satisfies $p\le K^4$. It does not disprove a bound
$K^{4+o(1)}$.

\section{Arithmetic frontier}

The exact order-three event now has several compatible forms:
\[
\begin{array}{ccl}
\text{all ratios:}
 &I\ge3
 &\iff \pi_0=\pi_1=\pi_2=0,\\[2mm]
 & &\iff \operatorname{corank}\operatorname{Sylv}(P,Q)\ge3,\\[2mm]
 & &\iff c_0=c_1=c_2=0,\\[3mm]
\text{inverse ratio:}
 &I_{-1}\ge3
 &\iff \deg\gcd(E,O)\ge2,\\[2mm]
 & &\iff \rho(0)=\rho'(0)=0,\\[2mm]
 & &\iff \exists\{x_1,x_2,x_3,x_4\}:e_1=e_3=0,\\[2mm]
 & &\iff \exists\text{ disjoint }e,f:\ \mathfrak q\supset(w_e,w_f).
\end{array}
\]
The characteristic-zero affine-intercept theorem ensures that the family does
not lie on the order-three stratum before reduction. Finite exceptional primes
are therefore genuine rank-drop primes of a structured cyclotomic system.

The remaining arithmetic theorem is not a height bound for one resultant. It
is common-divisor control for the small canonical systems above. A sufficient
global form remains, for example,
\[
P^+_{\rm split}\!\left(\prod_{[r]}\Pi_{K,r}^{(3)}\right)
 \le K^{4+o(1)}.
\]
A fixed-ratio logarithmic support budget $O(K^3)$ is already known, but it does
not imply this largest-prime statement. The $K=89$ certificate shows that a
strict finite $K^4$ ceiling is already too strong, so a successful theorem must
allow at least constant-factor excursions at the exponent-four scale.

The Fitting formulation points to concrete attacks: identities among
neighboring Sylvester minors (Desnanot--Jacobi/Dodgson), subresultant PRS or
Habicht identities, Bezoutian and reduced-resultant formulations, and modular
Smith/Fitting invariants. The local inverse formulation points instead to gcd
and norm control for two explicit sums of squares of cyclotomic units. The goal
is to control common split prime support, not individual norms.

\section{Verification boundary}

The repository maintains exact symbolic and finite-field regressions for every
coordinate system above. Principal-subresultant, inverse transversality and
reduced-edge gates run together with independent C++20 checks, Lean 4 and TeX
builds. Finite computations verify implementations and explicit anchors; they
are not the proof of the universal statements.

Representative valuation controls are
\[
\Pi'_{13,12}=5^3,
\qquad
\Lambda'_{13,12}=5^2,
\]
and
\[
\Theta'_{23}=47^2\cdot137,
\qquad
\Xi'_{23}=47\cdot137.
\]
Thus equality of split prime support must never be silently upgraded to
integer or valuation equality. Likewise,
\[
p\mid\gcd(N(w_e),N(w_f))
\]
does not imply that one prime ideal above $p$ divides both elements; the
same-embedding check is mandatory.

None of these row-support theorems proves microscopic BGS/GOE/COE universality.
Even a completed pointwise fourth-moment theorem leaves the independent
Heisenberg-scale joint-trace comparison and high-frequency tail problems.

\begin{thebibliography}{9}

\bibitem{BrownTraub}
W.~S.~Brown and J.~F.~Traub,
\emph{On Euclid's Algorithm and the Theory of Subresultants},
J. ACM \textbf{18} (1971), 505--514,
doi:10.1145/321662.321665.

\bibitem{Bass}
H.~Bass,
\emph{Generators and Relations for Cyclotomic Units},
Nagoya Math. J. \textbf{27} (1966), 401--407,
doi:10.1017/S0027763000026258.

\bibitem{Fu}
H.~Fu,
\emph{Projective Equivalence for the Roots of Unity},
J. Algebraic Combin. \textbf{56} (2022), 823--871,
arXiv:1905.01510.

\end{thebibliography}

\end{document}